\documentclass[11pt]{article}

\usepackage[margin=1in]{geometry}
\usepackage{amsmath,amssymb}
\usepackage{braket}
\usepackage{booktabs}
\usepackage{tabularx}
\usepackage{enumitem}
\usepackage{xcolor}
\usepackage[colorlinks=true, linkcolor=blue, citecolor=blue, urlcolor=blue]{hyperref}
\usepackage{listings}

% ── Listings (Julia) ────────────────────────────────────────────────────────
\definecolor{codebg}{rgb}{0.97,0.97,0.97}
\definecolor{codekey}{rgb}{0.0,0.0,0.6}
\definecolor{codestr}{rgb}{0.0,0.4,0.0}
\definecolor{codecom}{rgb}{0.4,0.4,0.4}

\lstdefinelanguage{Julia}{
  morekeywords={function,end,struct,abstract,type,if,else,elseif,for,in,while,
                return,using,import,include,module,const,where,mutable,do,let,
                begin,Vector,NTuple,Int,Float64,ComplexF64,T,Real,Number,Symbol,
                Bool,Array,Matrix,Tuple,Dict,nothing,true,false,Val},
  sensitive=true,
  morecomment=[l]{\#},
  morestring=[b]",
}
\lstset{
  language=Julia,
  basicstyle=\ttfamily\footnotesize,
  backgroundcolor=\color{codebg},
  keywordstyle=\color{codekey}\bfseries,
  stringstyle=\color{codestr},
  commentstyle=\color{codecom}\itshape,
  breaklines=true,
  showstringspaces=false,
  frame=single,
  framerule=0pt,
  xleftmargin=0.5em,
  xrightmargin=0.5em,
  tabsize=2,
}

\title{\textbf{Porting Plan: U(1)-Symmetric Tensor Networks for \texttt{Quench\_dyn}}\\
\large File-by-file specification for charge-conserving DMRG and TDVP\\ on the 1D XXZ chain}
\author{Nishan Ranabhat}
\date{\today}

\begin{document}
\maketitle

\begin{abstract}
This document specifies, in detail, the porting of the U(1)-symmetric
tensor-network machinery already implemented in \texttt{KagomeDSF/U1/} to the
1D XXZ project \texttt{Quench\_dyn}.  The motivation is the quench-dynamics
study laid out in \texttt{docs/manuscript/XXZ\_phase\_ordering.tex}: reaching
the regime where the even--odd splitting $\Delta_{\rm eo}(N) \sim e^{-N/\xi}$
becomes smaller than the random-field-induced mixing requires larger $N$ and
larger $\Delta$ than dense ED or dense MPS can access.  Exploiting total-$S^z$
conservation gives a $3\!-\!10\times$ reduction in memory and runtime at fixed
$\chi$, equivalently a larger effective $\chi$ at fixed cost.  No new
algorithms are needed; the existing U(1) layer is generic and the only
project-specific code is a $\sim 300$~LOC XXZ MPO builder.  The plan below
identifies (i) which files are copied verbatim, (ii) what new native code is
written, (iii) the test plan against the dense baseline and ED, and
(iv) the phased order of work.  Implementation is deferred to a separate
working directory.
\end{abstract}

\tableofcontents

% ===========================================================================
\section{Motivation and scope}
% ===========================================================================

The companion manuscript shows that the late-time N\'eel order following the
quench
\[
  H_i = H_{\rm XXZ}(\Delta_i) - \sum_j h_j\,S_j^z
  \;\longrightarrow\;
  H_f = H_{\rm XXZ}(\Delta_f)
\]
is parity-controlled: the local magnetization $\langle S_i^z(t)\rangle$ is
purely a cross-parity matrix element and scales as $\mathcal{O}(h)$ as long as
the random-field mixing $|\bra{\text{even}} V \ket{\text{odd}}|$ is small
compared with the even--odd splitting $\Delta_{\rm eo}(N) = 2|t|$.  In the
deep ordered regime the tunneling between the two N\'eel branches behaves as
$t \sim J(1/\Delta)^{cN} \sim e^{-N/\xi}$.  Reaching the crossover
$|v| \gtrsim \Delta_{\rm eo}(N)$ where spontaneous symmetry breaking should
become visible therefore requires \emph{both} large $\Delta$ and large $N$.
ED is restricted to $N \lesssim 20$; dense MPS is bottlenecked by the bond
dimension $\chi$ required to resolve the long-range entanglement built up by
the post-quench dynamics.

Both $H_i$ and $H_f$ commute with the total magnetization
$S^z_{\rm tot} = \sum_j S_j^z$, including the disorder term
$V = -\sum_j h_j S_j^z$, which is diagonal in $S^z$.  The entire protocol
therefore stays inside the $m=0$ sector (for even $N$), exactly the sector
that contains the two N\'eel branches $\ket{A}=\ket{\uparrow\downarrow\dots}$
and $\ket{B}=\ket{\downarrow\uparrow\dots}$.  Restricting MPS/MPO tensors to
the block-sparse U(1) representation is a strict gain with no physics lost.
The remaining $\mathbb{Z}_2$ parity inside the $m=0$ sector is left intact as
an open symmetry one can use as a diagnostic at the analysis stage.

\paragraph{Why a port rather than a rewrite.}
The U(1) layer already exists at \texttt{/home/nishan/UVA\_postdoc/Research/KagomeDSF/U1/}.
It was developed for the sawtooth/Kagome strip, which is genuinely harder
than 1D XXZ (supersites of dimension $4$, MPO bond dimension up to $\chi_{\rm MPO}=5$
with mixed charge content, two-sublattice geometry).  Of $\sim 6600$ lines of
Julia, roughly $5400$ are generic block-sparse machinery and only
$\sim 500$ are sawtooth-specific.  The XXZ chain is a strict subset of that
geometry: spin-$\tfrac12$ physical sites and a $\chi_{\rm MPO}=5$ bond with the
same five FSM states (Id-start, $S^+$-pending, $S^-$-pending, $S^z$-pending,
Id-finish) carrying charges $(0,+2,-2,0,0)$.  The dense XXZ MPO builder
already in \texttt{Quench\_dyn} (\texttt{src/Builders/xxz\_builder.jl}) uses
exactly these bond states with the upper-triangular convention, so the U(1)
projection is mechanical.

\paragraph{Out of scope.}
This document plans the U(1) port only.  We do \emph{not} remove or modify
the existing dense code under \texttt{src/}; the dense path remains the
reference implementation for validation and the small-$N$ workflow.  We do
\emph{not} change algorithm choices --- the existing DMRG, two-site TDVP, and
one-site TDVP carry over with the same APIs.  We do \emph{not} add new
symmetries (parity, $\mathbb{Z}_n$) at this stage; the U(1) layer is already
written generically over an \texttt{AbstractSymmetry} type and can be
extended later.

% ===========================================================================
\section{Symmetry analysis specific to our Hamiltonian}
% ===========================================================================

\subsection{Conservation of total $S^z$}
The XXZ Hamiltonian
\[
  H_{\rm XXZ} = J\sum_i \bigl(S_i^x S_{i+1}^x + S_i^y S_{i+1}^y\bigr)
              + \Delta\sum_i S_i^z S_{i+1}^z
\]
has $XY$ part $\tfrac{J}{2}\sum_i (S_i^+ S_{i+1}^- + S_i^- S_{i+1}^+)$ which
flips one spin up and a neighbouring spin down --- total $S^z$ unchanged.
The Ising part is diagonal in $S^z$, hence commutes trivially.  The disorder
term $V = -\sum_j h_j S_j^z$ is also diagonal in $S^z$.  Therefore
\[
  [H_i, S^z_{\rm tot}] = [H_f, S^z_{\rm tot}] = 0,
\]
and the time evolution remains in a fixed magnetization sector throughout
both the imaginary-time preparation (DMRG) and the real-time quench (TDVP).

\subsection{Charge convention}
Throughout we use the integer charge
\[
  q \;=\; 2S^z \;\in\; \mathbb{Z}.
\]
For spin-$\tfrac12$, $\ket{\uparrow} \leftrightarrow q=+1$ and
$\ket{\downarrow} \leftrightarrow q=-1$.  The action of $S^+$ raises $q$ by
$+2$; $S^-$ lowers by $-2$; $S^z$ preserves $q$.  This is the TeNPy
convention adopted by the existing U(1) layer
(\texttt{KagomeDSF/U1/src/Symmetry/chargeinfo.jl}).

\subsection{Bond charges on the XXZ MPO}
The dense XXZ MPO in \texttt{Quench\_dyn} uses a $\chi_{\rm MPO}=5$
upper-triangular bond with the FSM states already documented inline at
\texttt{src/Builders/xxz\_builder.jl:47--63}:
\begin{center}
\begin{tabular}{rlc}
\toprule
qindex & FSM state & bond charge $q$ \\
\midrule
1 & Id-start \;(no coupling in flight)        & $0$ \\
2 & $S^+$-placed \;(awaiting $S^-$ to the right) & $+2$ \\
3 & $S^-$-placed \;(awaiting $S^+$ to the right) & $-2$ \\
4 & $S^z$-placed \;(awaiting $S^z$ to the right) & $0$ \\
5 & Id-finish \;(coupling completed)             & $0$ \\
\bottomrule
\end{tabular}
\end{center}
With the U(1) layer's per-tensor rule on the MPO
$W[w_L, w_R, p, p^*]$,
\[
  \zeta_{w_L}\,q_{w_L}
  + \zeta_{w_R}\,q_{w_R}
  + \zeta_{p}\,q_{p}
  + \zeta_{p^*}\,q_{p^*}
  \;=\; 0,
  \qquad \zeta = (+1,-1,+1,-1),
\]
each non-zero entry of the dense $W$ already constructed by
\texttt{build\_xxz\_mpo} satisfies the rule.  Two spot checks:
\begin{itemize}[leftmargin=2em]
  \item $W[1,2,\uparrow,\downarrow] = (J/2)\,S^+$:
        $0 - 2 + 1 - (-1) = 0$ \checkmark
  \item $W[2,5,\downarrow,\uparrow] = S^-$:
        $+2 - 0 + (-1) - 1 = 0$ \checkmark
\end{itemize}
So the U(1) projection of the existing dense MPO is well-defined and the
``dense-then-\texttt{from\_dense}'' construction used in
\texttt{KagomeDSF/U1/src/Models/sawtooth.jl} carries over verbatim.

\subsection{The $m=0$ N\'eel branches and what U(1) does \emph{not} resolve}
Both N\'eel branches $\ket{A}$ and $\ket{B}$ have $S^z_{\rm tot}=0$ for even $N$;
they live in the same U(1) block.  Parity $P: S^z_j \mapsto -S^z_j$ maps $A
\leftrightarrow B$ \emph{within} the $m=0$ sector but maps $m \leftrightarrow
-m$ between sectors --- so parity does not commute with $S^z_{\rm tot}$
globally.  Inside $m=0$ alone, parity acts as a residual $\mathbb{Z}_2$.  We
do not enforce it as a symmetry in the tensor structure; the even/odd
splitting is what we are trying to \emph{resolve}, not symmetrize away.

% ===========================================================================
\section{Inventory of the existing U(1) layer}
% ===========================================================================

The U(1) package at \texttt{KagomeDSF/U1/} is self-contained: stdlib only
(\texttt{LinearAlgebra}, \texttt{Random}, \texttt{Printf}, \texttt{TOML})
plus \texttt{JLD2} for serialization.  It declares \texttt{TensorOperations}
in \texttt{Project.toml} but does not actually import it inside the U(1)
code.  The top-level layout (cf.\ \texttt{KagomeDSF/U1/src/U1.jl}):

\begin{center}
\renewcommand{\arraystretch}{1.05}
\begin{tabularx}{\linewidth}{l r X}
\toprule
Directory & LOC & Purpose \\
\midrule
\texttt{Symmetry/}  & $\sim 900$  & Charge metadata, leg-charge bookkeeping, leg pipes (fusion). \\
\texttt{Tensor/}    & $\sim 2000$ & \texttt{ChargedTensor} type, dense/charged adapters, linalg, contractions, SVD/QR. \\
\texttt{MPS/}       & $\sim 1050$ & Environments, canonicalization, generic builders, padding, IO. \\
\texttt{Algorithms/}& $\sim 1000$ & Effective Hamiltonians, Lanczos, Krylov-exponential, DMRG, TDVP (1- and 2-site). \\
\texttt{Models/}    & $\sim 500$  & Sawtooth-specific MPO and site builders. \emph{Not portable.} \\
\texttt{DSF/}       & $\sim 650$  & Spectral-function pipeline. \emph{Not needed for quench dynamics.} \\
\texttt{ED/}        & $\sim 600$  & Dense ED for sawtooth, validation only. \emph{Not needed.} \\
\texttt{Parallel/}  & --          & Threaded reimplementations of hot paths. \emph{Optional, post-MVP.} \\
\bottomrule
\end{tabularx}
\end{center}

\subsection{Key types and conventions to inherit}

\textbf{Charge metadata (\texttt{Symmetry/chargeinfo.jl}).}
\begin{lstlisting}
abstract type AbstractSymmetry end
struct U1Symmetry <: AbstractSymmetry end

struct ChargeInfo{S<:AbstractSymmetry}
    mod::Int      # 0 for U(1)
    name::Symbol  # :twice_Sz
end
\end{lstlisting}

\textbf{Leg charges (\texttt{Symmetry/legcharge.jl}).}
\begin{lstlisting}
struct LegCharge{S<:AbstractSymmetry}
    chinfo::ChargeInfo{S}
    charges::Vector{Int}    # one per qindex (block)
    slices::Vector{Int}     # block boundaries, length nblocks+1
    qconj::Int8             # +-1
end
\end{lstlisting}

\textbf{Block-sparse tensor (\texttt{Tensor/chargedtensor.jl}).}
\begin{lstlisting}
struct ChargedTensor{T, N, S<:AbstractSymmetry}
    legs::NTuple{N, LegCharge{S}}
    qtotal::Int
    qdata::Matrix{Int}        # (nblocks_stored, N) of qindices
    data::Vector{Array{T, N}} # the block contents
end
\end{lstlisting}

\textbf{Sign convention (the \emph{single} thing that matters most when porting):}
\begin{center}
\begin{tabular}{lll}
\toprule
Tensor & Legs & $\zeta$ \\
\midrule
MPS site & $(v_L, p, v_R)$        & $(+1, +1, -1)$ \\
MPO site & $(w_L, w_R, p, p^*)$   & $(+1, -1, +1, -1)$ \\
Left env & $(\bar v, w, v)$       & $(+1, -1, -1)$ \\
Right env& $(\bar v, w, v)$       & $(-1, +1, +1)$ \\
\bottomrule
\end{tabular}
\end{center}
with the per-tensor charge rule
$\sum_k \zeta_k\,q_{k,\,\text{qidx}[k]} \equiv Q\pmod{\text{mod}}$ for the
stored qindex tuples.

\subsection{Operations already implemented}

\begin{itemize}[leftmargin=2em]
\item \textbf{Adapters.} \texttt{from\_dense(A, legs, qtotal; tol)} projects a
      dense array onto its allowed-sector blocks, warning if discarded mass
      exceeds \texttt{tol}.  \texttt{to\_dense(t)} inflates back.
\item \textbf{Linalg.} \texttt{conj}, \texttt{norm}, \texttt{dot}, scalar
      ops, \texttt{axpby}; everything per-block.
\item \textbf{Contractions.} \texttt{permute}, block-sparse
      \texttt{tensordot}, \texttt{combine\_legs}/\texttt{split\_legs} via
      \texttt{LegPipe}.
\item \textbf{Decompositions.} \texttt{svd\_charged} with per-block
      $\sigma$-floor, global sort, relative cutoff, hard cap on $\chi$.
      \texttt{qr\_charged} for canonicalization.
\item \textbf{MPS ops.} \texttt{trivial\_left\_env}, \texttt{update\_left\_env},
      symmetric right-side variants, \texttt{build\_environments},
      \texttt{make\_canonical!}, orthogonality checks, \texttt{measure\_norm},
      \texttt{measure\_energy}, \texttt{pad\_mps\_charged} (charge-preserving).
\item \textbf{Solvers.} \texttt{TwoSiteEffectiveH}, \texttt{OneSiteEffectiveH},
      \texttt{ZeroSiteEffectiveH}, each with \texttt{\_apply};
      \texttt{LanczosSolver}, \texttt{KrylovExponential}, each acting
      directly on \texttt{ChargedTensor}.
\item \textbf{Algorithms.} \texttt{ChargedMPSState} container,
      \texttt{dmrg\_sweep!}, \texttt{tdvp\_sweep!} (2-site),
      \texttt{tdvp\_sweep\_one\_site!} (1-site).
\end{itemize}

% ===========================================================================
\section{Target architecture for the port}
% ===========================================================================

\subsection{Where the new code lives}
We propose creating a sibling sub-package under \texttt{Quench\_dyn/} that
mirrors the \texttt{KagomeDSF/U1/} layout:
\begin{verbatim}
Quench_dyn/
  src/                <-- existing dense code, untouched
  U1/                 <-- NEW sub-package, self-contained
    Project.toml
    src/
      QuenchDynU1.jl  <-- module top, includes everything below
      Symmetry/       <-- copy from KagomeDSF/U1/src/Symmetry/
      Tensor/         <-- copy from KagomeDSF/U1/src/Tensor/
      MPS/            <-- copy from KagomeDSF/U1/src/MPS/
      Algorithms/     <-- copy from KagomeDSF/U1/src/Algorithms/
      Models/
        xxz.jl        <-- NEW (~300 LOC), replaces sawtooth.jl
    test/
      ... (copied generic tests + new XXZ tests)
\end{verbatim}
This keeps the dense and U(1) code paths fully decoupled.  The existing
scripts in \texttt{scripts/} continue to work against \texttt{QuenchDyn}; new
scripts that want U(1) load \texttt{QuenchDynU1} instead.  Once the U(1)
path is stable we can decide whether to merge them under one umbrella module;
that is deferred.

\subsection{Why a sub-package and not a flat add}
Two reasons.  First, it mirrors how \texttt{KagomeDSF/U1/} was structured ---
proven to work, so we inherit the dependency hygiene for free.  Second, a
sub-package has its own \texttt{Project.toml} and tests, so the U(1) layer
can be precompiled and tested independently of the dense code.  This matches
your modular-architecture preference (no black-box pipelines, every
operation independent).

\subsection{Renaming considerations during the copy}
None.  Every type, function, and convention in the U(1) layer is XXZ-agnostic
already.  The module name (\texttt{U1} vs \texttt{QuenchDynU1}) is the only
thing that needs to change at the top of \texttt{src/QuenchDynU1.jl}, and
that is just an edit to the \texttt{module} keyword and the file header.

% ===========================================================================
\section{Phase A--D: portable layers (copy verbatim)}
% ===========================================================================

\subsection{Phase A --- Symmetry (3 files)}
\begin{center}
\begin{tabularx}{\linewidth}{l r X}
\toprule
File & LOC & Notes \\
\midrule
\texttt{Symmetry/chargeinfo.jl} & 159 & \texttt{AbstractSymmetry}, \texttt{NoSymmetry}, \texttt{U1Symmetry}, \texttt{ChargeInfo}, \texttt{make\_valid}.  Copy as-is. \\
\texttt{Symmetry/legcharge.jl}  & $\sim 340$ & \texttt{LegCharge}, \texttt{from\_qflat}, \texttt{from\_trivial}, \texttt{block\_size(s)}, \texttt{slice}, \texttt{charge}, \texttt{findblock}, \texttt{is\_contractible}.  Copy as-is. \\
\texttt{Symmetry/legpipe.jl}    & $\sim 400$ & \texttt{LegPipe}, \texttt{q\_map}, \texttt{strides}, helpers used by \texttt{combine\_legs}.  Copy as-is. \\
\bottomrule
\end{tabularx}
\end{center}

\subsection{Phase B --- Tensor (5 files)}
\begin{center}
\begin{tabularx}{\linewidth}{l r X}
\toprule
File & LOC & Notes \\
\midrule
\texttt{Tensor/chargedtensor.jl}  & $\sim 450$ & The \texttt{ChargedTensor\{T,N,S\}} type and unchecked/checked constructors. \\
\texttt{Tensor/adapters.jl}       & $\sim 250$ & \texttt{validate\_charges}, \texttt{from\_dense}, \texttt{to\_dense}. \\
\texttt{Tensor/linalg.jl}         & $\sim 200$ & \texttt{conj}, \texttt{copy}, \texttt{norm}, \texttt{dot}, scalar ops, \texttt{axpby}. \\
\texttt{Tensor/contraction.jl}    & $\sim 600$ & \texttt{permute}, \texttt{tensordot}, \texttt{combine\_legs}, \texttt{split\_legs}. \\
\texttt{Tensor/decomposition.jl}  & $\sim 500$ & \texttt{svd\_charged}, \texttt{qr\_charged}, truncation policy. \\
\bottomrule
\end{tabularx}
\end{center}

\subsection{Phase C --- MPS operations (4 files)}
\begin{center}
\begin{tabularx}{\linewidth}{l r X}
\toprule
File & LOC & Notes \\
\midrule
\texttt{MPS/tensorops.jl} & $\sim 450$ & \texttt{trivial\_left\_env}/\texttt{right\_env}, \texttt{update\_left\_env}, symmetric right, \texttt{build\_environments}, \texttt{move\_orthogonality\_*!}, \texttt{make\_canonical!}, \texttt{is\_left/right\_orthogonal}, \texttt{measure\_norm}, \texttt{measure\_energy}. \\
\texttt{MPS/builders.jl} & $\sim 350$ & Generic \texttt{spin\_half\_leg}, \texttt{trivial\_boundary\_leg}, \texttt{product\_state\_charged}, \texttt{random\_charged\_mps}, \texttt{rand\_charged\_tensor}.  Sawtooth-specific functions (\texttt{sawtooth\_supersite\_leg}, \texttt{sawtooth\_mpo\_bond\_leg}, \texttt{sawtooth\_random\_charged\_mps}, \texttt{sawtooth\_supersite\_operator}) should be deleted on copy --- they belong in \texttt{Models/}.\\
\texttt{MPS/io.jl}       & $\sim 150$ & \texttt{save\_charged\_mps}, \texttt{load\_charged\_mps}.  Copy as-is. \\
\texttt{MPS/padding.jl}  & $\sim 100$ & \texttt{pad\_mps\_charged}.  Charge-preserving padding; essential for the 1-site TDVP workflow. \\
\bottomrule
\end{tabularx}
\end{center}

\paragraph{One pruning during copy.}  \texttt{MPS/builders.jl} mixes generic
constructors with sawtooth supersite helpers.  The split is clean: anything
named \texttt{sawtooth\_*} goes; anything else is generic and stays.  The
generic \texttt{product\_state\_charged} takes a physical \texttt{LegCharge}
and a vector of basis indices --- exactly what we need to construct N\'eel
states for the XXZ chain (Section~\ref{sec:phase-f}).

\subsection{Phase D --- Algorithms (2 files)}
\begin{center}
\begin{tabularx}{\linewidth}{l r X}
\toprule
File & LOC & Notes \\
\midrule
\texttt{Algorithms/solvers.jl}    & $\sim 450$ & \texttt{EffectiveHamiltonian} hierarchy (zero-/one-/two-site), \texttt{\_apply}, \texttt{LanczosSolver} (\texttt{\_solve}), \texttt{KrylovExponential} (\texttt{\_evolve}). \\
\texttt{Algorithms/algorithms.jl} & $\sim 550$ & \texttt{ChargedMPSState}, \texttt{dmrg\_sweep!}, \texttt{tdvp\_sweep!} (2-site), \texttt{tdvp\_sweep\_one\_site!}. \\
\bottomrule
\end{tabularx}
\end{center}

\paragraph{API mapping (dense $\to$ U(1)).}
\begin{center}
\begin{tabularx}{\linewidth}{l l X}
\toprule
Dense (\texttt{QuenchDyn}) & U(1) (\texttt{QuenchDynU1}) & Differences \\
\midrule
\texttt{MPS\{T\}}      & \texttt{Vector\{ChargedTensor\{T,3,U1Symmetry\}\}} & No wrapper struct in U(1); a bare \texttt{Vector}. \\
\texttt{MPO\{T\}}      & \texttt{Vector\{ChargedTensor\{T,4,U1Symmetry\}\}} & Same. \\
\texttt{Environment\{T\}} & \texttt{Vector\{Union\{Nothing, ChargedTensor\{T,3,U1Symmetry\}\}\}} & Same. \\
\texttt{MPSState}      & \texttt{ChargedMPSState\{Tmps,Tmpo,Tenv,S\}} & Carries explicit eltypes; same role. \\
\texttt{DMRGOptions}   & passed as \texttt{(chi\_max, cutoff)} kwargs & U(1) inlines the options into \texttt{dmrg\_sweep!}. \\
\texttt{TDVPOptions}   & passed as \texttt{(dt, chi\_max, cutoff)} kwargs & Same comment. \\
\texttt{dmrg\_sweep}   & \texttt{dmrg\_sweep!} & In-place in U(1) (matches existing convention there). \\
\texttt{tdvp\_sweep}   & \texttt{tdvp\_sweep!} & In-place; otherwise identical sweep mechanics. \\
\texttt{tdvp\_sweep\_one\_site} & \texttt{tdvp\_sweep\_one\_site!} & Same. \\
\bottomrule
\end{tabularx}
\end{center}
The \texttt{!} naming difference (mutating in U(1) vs returning new in dense)
is a real semantic distinction; we keep U(1)'s convention as it is and
\emph{do not} rewrite the dense routines.

% ===========================================================================
\section{Phase E --- Native: the XXZ charged MPO}
\label{sec:phase-e}
% ===========================================================================

This is the only piece of model-specific code we write.  The strategy
mirrors \texttt{KagomeDSF/U1/src/Models/sawtooth.jl}: build the dense $W$
tensors with the existing routine, declare the leg-charge structure, and
project with \texttt{from\_dense}.

\subsection{File: \texttt{U1/src/Models/xxz.jl}}

\paragraph{Public functions.}
\begin{lstlisting}
"""
    build_xxz_mpo_charged(N, J, Delta, h::Vector{Float64};
                          T::Type = Float64) -> Vector{ChargedTensor{T,4,U1Symmetry}}

U(1)-charge-conserving XXZ MPO with chi_MPO = 5 and per-site longitudinal
field h_i. Charge convention: q = 2 S^z. Bond states and charges:
  [Id-start, S+-pending, S--pending, Sz-pending, Id-finish]
  charges = [0, +2, -2, 0, 0]
"""
function build_xxz_mpo_charged(N::Int, J::Float64, Delta::Float64,
                               h::Vector{Float64};
                               T::Type = Float64)
end

# Convenience overloads mirroring the dense API:
build_xxz_mpo_charged(N, J, Delta, h_uniform::Float64; T=Float64)
build_xxz_mpo_charged(N, J, Delta;                       T=Float64)
\end{lstlisting}

\paragraph{Leg constructors.}
\begin{lstlisting}
"""Physical spin-1/2 leg with charges (+1, -1). Bunched (no degeneracy)."""
xxz_phys_leg(ci::ChargeInfo{U1Symmetry}; qconj=+1) =
    LegCharge(ci, [(+1, 1), (-1, 1)], qconj)

"""chi_MPO = 5 XXZ bond leg in unbunched form: five qindices, each dim 1,
charges (0, +2, -2, 0, 0) corresponding to the FSM states above."""
xxz_mpo_bond_leg(ci::ChargeInfo{U1Symmetry}; qconj=+1) =
    LegCharge(ci, [0, +2, -2, 0, 0], [1, 2, 3, 4, 5, 6], qconj)
\end{lstlisting}
The bunched physical leg \texttt{xxz\_phys\_leg} is essentially an alias of
the existing generic \texttt{spin\_half\_leg} from \texttt{MPS/builders.jl}.
Defining it locally inside \texttt{Models/xxz.jl} keeps the XXZ vocabulary
self-contained and avoids reaching into the generic namespace from physics
code.  Bunching the three q=0 bond states (states 1, 4, 5) is a $\sim 1.5\times$
speedup on $W$-block contractions; the sawtooth blueprint explicitly leaves
this as a post-MVP refinement, so we follow the same policy and keep the
five states unbunched in the initial port.

\paragraph{MPO assembly (pseudocode).}
\begin{lstlisting}
function build_xxz_mpo_charged(N, J, Delta, h::Vector{Float64};
                               T::Type = Float64)
    @assert length(h) == N
    @assert N >= 2

    # 1. Dense reference: reuse Quench_dyn's existing builder.
    W_dense = QuenchDyn.build_xxz_mpo(N, J, Delta, h; T = T).tensors
    # W_dense is a Vector of 4-arrays sized (chi_l, chi_r, d, d), with
    # boundary tensors already sliced to (1, 5, d, d) and (5, 1, d, d).

    # 2. Leg structure.
    ci             = ChargeInfo{U1Symmetry}(0, :twice_Sz)
    phys           = xxz_phys_leg(ci)
    bond_boundary  = trivial_boundary_leg(ci)    # dim-1, q = 0
    bond_bulk      = xxz_mpo_bond_leg(ci)

    legs_first = (bond_boundary, conj(bond_bulk), phys, conj(phys))
    legs_bulk  = (bond_bulk,     conj(bond_bulk), phys, conj(phys))
    legs_last  = (bond_bulk,     conj(bond_boundary), phys, conj(phys))

    # 3. Project to U(1) blocks. Every non-zero entry of the dense W
    # satisfies the charge rule, so discarded mass should be ~0
    # (warn threshold = 1e-12).
    mpo = Vector{ChargedTensor{T, 4, U1Symmetry}}(undef, N)
    mpo[1] = from_dense(W_dense[1], legs_first, 0)
    for i in 2:(N - 1)
        mpo[i] = from_dense(W_dense[i], legs_bulk, 0)
    end
    mpo[N] = from_dense(W_dense[N], legs_last, 0)
    return mpo
end
\end{lstlisting}

\paragraph{Per-site vs shared tensors.}  Site-dependent $h_i$ means we
cannot share the same \texttt{bulk} tensor across sites (as the sawtooth
builder does); each site needs its own \texttt{from\_dense} call.  This is
fine --- the dense XXZ builder already constructs $N$ distinct tensors.  For
the special case $h_i \equiv h_0$ (uniform field) one could share a single
bulk \texttt{ChargedTensor} reference, but the memory saving is negligible
($\chi_{\rm MPO}=5$, $d=2$ means each $W$ is at most $5 \cdot 5 \cdot 4 = 100$
floats).  We do not optimize this.

\paragraph{Cross-dependency note.}  The pseudocode above shows \texttt{QuenchDynU1.build\_xxz\_mpo\_charged}
calling \texttt{QuenchDyn.build\_xxz\_mpo} to reuse the dense $W$ tensors.
This introduces a one-way dependency \texttt{QuenchDynU1 $\to$ QuenchDyn}.
If we want \texttt{QuenchDynU1} to be standalone (mirroring how
\texttt{KagomeDSF/U1} is standalone) we should instead inline the
dense $W$ construction --- $\sim 50$ LOC --- into \texttt{Models/xxz.jl}.
Recommendation: inline.  It removes the cross-dep, the code is small, and the
two builders can diverge in future without coupling.

% ===========================================================================
\section{Phase F --- Native: N\'eel state and quench-protocol initialisers}
\label{sec:phase-f}
% ===========================================================================

\paragraph{N\'eel product state in the $m=0$ sector.}
\begin{lstlisting}
"""
    xxz_neel_state(::Type{T}, N::Int; start_up::Bool = true)
        -> Vector{ChargedTensor{T, 3, U1Symmetry}}

The charge-conserving Neel product state in the m = 0 sector of the
spin-1/2 chain. start_up = true gives |up dn up dn ...>, false gives
|dn up dn up ...>. Requires even N.
"""
function xxz_neel_state(::Type{T}, N::Int; start_up::Bool = true) where T
    iseven(N) || throw(ArgumentError(
        "xxz_neel_state needs even N for m = 0; got N = $N."))
    ci   = ChargeInfo{U1Symmetry}(0, :twice_Sz)
    phys = xxz_phys_leg(ci)
    # phys basis indices: 1 = up (q=+1), 2 = dn (q=-1)
    indices = start_up ? [iseven(i) ? 2 : 1 for i in 1:N] :
                         [iseven(i) ? 1 : 2 for i in 1:N]
    return product_state_charged(T, phys, indices)
end
\end{lstlisting}
The two N\'eel branches $\ket{A}$ and $\ket{B}$ in the manuscript correspond
to \texttt{start\_up = true} and \texttt{false} respectively.

\paragraph{Random charged MPS in a target sector.}
We expose a thin wrapper around the existing generic constructor:
\begin{lstlisting}
"""
    xxz_random_charged_mps(::Type{T}, N::Int, chi::Int;
                           qtotal::Int = 0) -> Vector{ChargedTensor{T,3,U1Symmetry}}

Random charge-conserving spin-1/2 MPS with total charge qtotal (q = 2*S^z_tot).
Used for initial DMRG warm starts and as a richer-than-product seed for
benchmarking convergence.
"""
function xxz_random_charged_mps(::Type{T}, N::Int, chi::Int;
                                qtotal::Int = 0) where T
    ci   = ChargeInfo{U1Symmetry}(0, :twice_Sz)
    phys = xxz_phys_leg(ci)
    return random_charged_mps(T, phys, N, chi; qtotal = qtotal)
end
\end{lstlisting}

\paragraph{Local operators for measurements.}
\begin{lstlisting}
"""Single-site operator as a (2,2) dense matrix, embedded later
   via apply_local_operator_charged or measure_local_observable."""
xxz_local_op(symbol::Symbol; T = Float64) =
    Dict(:Z => T[0.5 0; 0 -0.5],
         :X => T[0 0.5; 0.5 0],
         :Sp => T[0 1; 0 0], :Sm => T[0 0; 1 0],
         :I => T[1 0; 0 1])[symbol]
\end{lstlisting}
Note that $S^z$ and $S^x$ are charge-diagonal (preserve $q$) and
charge-changing ($\Delta q = \pm 2$) respectively.  $\langle S^z(t)\rangle$
is the only observable the manuscript needs and is fully accessible in
sector-conserving form.

% ===========================================================================
\section{Module integration and exports}
% ===========================================================================

\paragraph{New module file: \texttt{U1/src/QuenchDynU1.jl}.}
\begin{lstlisting}
module QuenchDynU1

using LinearAlgebra
using Random
using JLD2

include("Symmetry/chargeinfo.jl")
include("Symmetry/legcharge.jl")
include("Symmetry/legpipe.jl")

include("Tensor/chargedtensor.jl")
include("Tensor/adapters.jl")
include("Tensor/linalg.jl")
include("Tensor/contraction.jl")
include("Tensor/decomposition.jl")

include("MPS/tensorops.jl")
include("MPS/builders.jl")
include("MPS/io.jl")
include("MPS/padding.jl")

include("Algorithms/solvers.jl")
include("Algorithms/algorithms.jl")

include("Models/xxz.jl")

# -- exports (the only XXZ-specific ones) -----------------------------
export xxz_phys_leg, xxz_mpo_bond_leg
export build_xxz_mpo_charged
export xxz_neel_state, xxz_random_charged_mps, xxz_local_op

# (plus the generic exports identical to U1.jl's, minus DSF/Models/ED ones)

end # module
\end{lstlisting}

\paragraph{\texttt{Project.toml}.}  Copy from \texttt{KagomeDSF/U1/Project.toml},
change \texttt{name} to \texttt{QuenchDynU1}, allocate a fresh UUID.  Drop
\texttt{TensorOperations} from \texttt{[deps]} since the U(1) code does not
actually use it.  Keep \texttt{JLD2}, \texttt{TOML} (for any future config
loaders).

% ===========================================================================
\section{Validation plan}
% ===========================================================================

The U(1) layer carries its own test suite at \texttt{KagomeDSF/U1/test/}.
We copy it (renaming where it references \texttt{U1} as a module name), then
\emph{add} XXZ-specific tests.

\subsection{Inherited generic tests (copy + rename)}
\begin{center}
\begin{tabularx}{\linewidth}{l X}
\toprule
File & What it checks \\
\midrule
\texttt{test\_chargeinfo.jl}    & \texttt{make\_valid}, U(1) vs $\mathbb{Z}_n$ edge cases. \\
\texttt{test\_legcharge.jl}     & Constructors, \texttt{is\_contractible}, \texttt{conj}, \texttt{slice}. \\
\texttt{test\_legpipe.jl}       & \texttt{q\_map} consistency, fusion-of-fusions. \\
\texttt{test\_chargedtensor.jl} & Class invariants, accessors, constructor validation. \\
\texttt{test\_adapters.jl}      & \texttt{to\_dense $\circ$ from\_dense} round-trip; off-sector projection mass. \\
\texttt{test\_linalg.jl}        & \texttt{norm}, \texttt{dot}, scalar mul, \texttt{axpby}. \\
\texttt{test\_contraction.jl}   & \texttt{tensordot} vs dense, \texttt{combine/split\_legs} round-trip. \\
\texttt{test\_decomposition.jl} & \texttt{svd\_charged} vs dense SVD, truncation policy. \\
\texttt{test\_tensorops.jl}     & Env build / update consistency; orthogonality after \texttt{make\_canonical!}. \\
\texttt{test\_solvers.jl}       & Lanczos ground state vs dense; Krylov-exp vs dense \texttt{exp(-iHt)}. \\
\texttt{test\_algorithms.jl}    & End-to-end DMRG/TDVP on a small generic Heisenberg-like Hamiltonian. \\
\texttt{test\_builders.jl}      & \texttt{product\_state\_charged}, \texttt{random\_charged\_mps} (after pruning sawtooth bits). \\
\texttt{test\_io.jl}            & JLD2 round-trip. \\
\bottomrule
\end{tabularx}
\end{center}

\subsection{New XXZ-specific tests}

\paragraph{\texttt{test\_xxz\_mpo\_charged.jl}.}
\begin{enumerate}[leftmargin=2em]
  \item For $N \in \{4, 6, 8\}$, $\Delta \in \{0.5, 1.0, 2.0\}$,
        random $h_j$ with $|h_j| \le 0.1$, and uniform $h$:
        verify \texttt{to\_dense(build\_xxz\_mpo\_charged(...)) == build\_xxz\_mpo(...).tensors[i]}
        block by block, to tolerance $10^{-12}$.
  \item Verify \texttt{validate\_charges} passes on every site.
  \item Verify the discarded mass returned by \texttt{from\_dense} is $< 10^{-13}$.
\end{enumerate}

\paragraph{\texttt{test\_xxz\_dmrg\_charged.jl}.}
\begin{enumerate}[leftmargin=2em]
  \item For $N=8$, several $\Delta$:
        run \texttt{dmrg\_sweep!} starting from \texttt{xxz\_neel\_state}.
        Compare ground-state energy against:
        \begin{itemize}
          \item dense DMRG (\texttt{QuenchDyn.dmrg\_sweep}), and
          \item ED (\texttt{build\_xxz\_hamiltonian} + \texttt{diagonalize}),
        \end{itemize}
        to $10^{-8}$ at $\chi=32$.
  \item Verify total charge is conserved exactly across all sweeps
        (\texttt{qtotal == 0} on the active site).
\end{enumerate}

\paragraph{\texttt{test\_xxz\_tdvp\_charged.jl}.}
\begin{enumerate}[leftmargin=2em]
  \item Trotter quench on $N=8$: prepare ground state of $H_i(\Delta_i=0.5, h)$,
        evolve under $H_f(\Delta_f=2.0)$ for $t \in [0, 5]$ at $dt=0.05$.
        Compare $\langle S_i^z(t)\rangle$ profiles vs ED (\texttt{ed\_evolve\_and\_measure}),
        to $10^{-6}$ at $\chi=64$.
  \item Compare against dense TDVP (\texttt{QuenchDyn.tdvp\_sweep}) at same parameters;
        verify agreement to $10^{-8}$.
  \item Check norm preservation in 1-site TDVP (should be machine precision).
  \item Check energy drift in 2-site TDVP (should match dense baseline).
\end{enumerate}

\paragraph{\texttt{test\_quench\_protocol\_charged.jl}.}
\begin{enumerate}[leftmargin=2em]
  \item Reproduce the manuscript's Fig.~1/2 setup on $N=12$ (still ED-comparable):
        DMRG with $h_j$ disorder, quench to clean $H_f$, evolve, plot
        $\langle S_i^z(t)\rangle$ and compare ED-style profile.
        This is a smoke test for the full protocol, not a unit test.
\end{enumerate}

\subsection{Drift benchmarks}
Mirror the existing \texttt{tests/benchmark\_tdvp\_drift.jl} for the U(1)
path.  We expect (i) lower memory at the same $\chi$ and (ii) the same
qualitative drift behaviour --- the U(1) layer does not change the
truncation policy, so 2-site TDVP still loses norm via SVD and 1-site TDVP
still preserves it.

% ===========================================================================
\section{Open questions and risks}
% ===========================================================================

\subsection{MPO convention drift}
The XXZ MPO in \texttt{Quench\_dyn} uses the upper-triangular FSM
convention; the sawtooth MPO in \texttt{KagomeDSF/U1} also uses an
upper-triangular form, but the dense reference in \texttt{KagomeDSF}'s
non-U(1) builder is lower-triangular.  This was a real bug source there
(\texttt{docs/mpo\_convention\_fix.md} discusses it).  Our path is safer:
we only invoke the U(1) projection on tensors produced by our own
\texttt{build\_xxz\_mpo}, whose convention is documented and consistent.
\textbf{Action:} during Phase E, the very first thing to check is that
\texttt{from\_dense} returns the expected non-zero qindex tuples on a
single bulk site (\texttt{nblocks\_stored == 9} for the XXZ $W$:
1 identity-diagonal entry on each of the three q=0 bond pairs, plus
$\{(2,5),(3,5),(4,5)\}$ on the right column, plus
$\{(1,2),(1,3),(1,4)\}$ in the first row).

\subsection{Padding compatibility for 1-site TDVP}
The dense 1-site routines in \texttt{Quench\_dyn} cannot grow $\chi$ and
require padding from a 2-site warm start.  In the U(1) case, padding must
preserve charge sectors --- the existing \texttt{pad\_mps\_charged} in
\texttt{KagomeDSF/U1/src/MPS/padding.jl} already does this.  \textbf{Action:}
test \texttt{xxz\_neel\_state $\to$ pad\_mps\_charged $\to$ 1-site DMRG}
end-to-end as part of Phase F.

\subsection{Numerical type uniformity}
The U(1) layer parameterises tensors over the element type \texttt{T}.
For real-time TDVP we will need \texttt{T = ComplexF64}.  DMRG can use
\texttt{Float64}.  The transition from a real DMRG ground state to complex
TDVP evolution happens by widening the type when we build the post-quench
state --- the dense code handles this implicitly; in the U(1) port we may
need an explicit \texttt{convert\_eltype} helper.  \textbf{Action:} verify
on a small case during Phase D--E.

\subsection{Performance expectations}
For 1D spin-$\tfrac12$ XXZ in the $m=0$ sector, the typical U(1) speedup
factor over dense MPS is $3\!-\!5\times$ in both memory and per-sweep time
at fixed $\chi$.  At fixed cost, that translates to a factor $\sim\sqrt{3}$
to $\sqrt{5}$ in effective $\chi$ --- enough to reach noticeably larger
$N$ in the ordered regime, but not orders of magnitude.  \textbf{Reality
check:} if the symmetry-broken regime in the manuscript demands
$\chi \sim 10^3\!-\!10^4$, the U(1) port alone may not be sufficient; we
would then also need block-aware threading
(\texttt{Parallel/} sub-package, currently out of scope) or a different
algorithm (controlled-bond-expansion TDVP, basis transformation).  This is
a decision for after we have measured concrete performance.

\subsection{Bunching the $q=0$ bond states}
The five XXZ MPO bond states include three with $q=0$ (Id-start, Sz-pending,
Id-finish).  They could be bunched into a single qindex of dimension 3,
halving the number of off-diagonal blocks the contractor walks through.  The
sawtooth blueprint defers this to ``post-MVP''.  We follow the same policy:
get correctness first with five unbunched qindices, profile, then bunch if
the speedup is meaningful.

% ===========================================================================
\section{Phased schedule}
% ===========================================================================

\begin{center}
\renewcommand{\arraystretch}{1.15}
\begin{tabularx}{\linewidth}{c X l}
\toprule
Phase & Deliverable & Estimate \\
\midrule
0 & Create \texttt{Quench\_dyn/U1/} skeleton, \texttt{Project.toml}, empty \texttt{QuenchDynU1.jl}. & 0.5 d \\
A & Copy \texttt{Symmetry/} (3 files).  Run \texttt{test\_chargeinfo.jl}, \texttt{test\_legcharge.jl}, \texttt{test\_legpipe.jl}. & 0.5 d \\
B & Copy \texttt{Tensor/} (5 files).  Run \texttt{test\_chargedtensor.jl}, \texttt{test\_adapters.jl}, \texttt{test\_linalg.jl}, \texttt{test\_contraction.jl}, \texttt{test\_decomposition.jl}. & 1 d \\
C & Copy \texttt{MPS/} (4 files), prune sawtooth helpers from \texttt{builders.jl}.  Run \texttt{test\_tensorops.jl}, \texttt{test\_builders.jl}, \texttt{test\_io.jl}. & 1 d \\
D & Copy \texttt{Algorithms/} (2 files).  Run \texttt{test\_solvers.jl}, \texttt{test\_algorithms.jl} on a generic Heisenberg test. & 1 d \\
E & Write \texttt{Models/xxz.jl}, \texttt{build\_xxz\_mpo\_charged}, leg helpers. Add \texttt{test\_xxz\_mpo\_charged.jl}. & 1--1.5 d \\
F & Add \texttt{xxz\_neel\_state}, \texttt{xxz\_random\_charged\_mps}, local-op helper.  Add \texttt{test\_xxz\_dmrg\_charged.jl}.  Confirm DMRG ground state vs dense and ED on $N=8$. & 1 d \\
G & Add \texttt{test\_xxz\_tdvp\_charged.jl}: 2-site TDVP quench vs dense and ED on $N=8$ and $N=12$. & 1--1.5 d \\
H & Padding $+$ 1-site TDVP path: \texttt{pad\_mps\_charged} round-trip, then 1-site TDVP smoke test. & 0.5 d \\
I & End-to-end protocol script reproducing manuscript Fig.~1/2 setup on $N=16$ (out of ED reach) at moderate $\chi$.  Visual comparison of $\langle S_i^z(t)\rangle$ vs dense TDVP. & 1 d \\
J & Drift benchmark: \texttt{benchmark\_tdvp\_drift\_charged.jl} mirroring the dense one. & 0.5 d \\
\midrule
& \textbf{Total} & $\sim 9\!-\!11$ working days \\
\bottomrule
\end{tabularx}
\end{center}

The phasing is intentionally bottom-up: every layer is tested independently
before the next is added.  Phases A--D consist almost entirely of copy and
test execution; Phase E is the only design work (the XXZ MPO charged
builder, $\sim 100$ LOC of actual code plus $\sim 200$ LOC of tests and
docstrings).  Phases F--J are integration and validation.

% ===========================================================================
\section{What this plan does \emph{not} commit to}
% ===========================================================================

\begin{itemize}[leftmargin=2em]
  \item Replacing the dense code under \texttt{src/}.  The dense path stays
        as the small-$N$ reference and is what the existing tests and scripts
        consume.  Both paths coexist.
  \item Enforcing the residual $\mathbb{Z}_2$ parity in the m=0 sector.  As
        argued in the manuscript, parity is what we want to \emph{resolve};
        baking it into the tensor structure would defeat the purpose.
  \item Multi-threaded execution (\texttt{Parallel/} sub-package in the
        sawtooth code).  Deferred until the serial U(1) path proves
        insufficient.
  \item Bunching the $q=0$ MPO bond states.  Deferred until profiling
        warrants it.
  \item Spectral-function / DSF infrastructure.  Not needed for the quench
        manuscript; the \texttt{DSF/} sub-package is left in
        \texttt{KagomeDSF/U1/}.
\end{itemize}

% ===========================================================================
\section{Summary}
% ===========================================================================

The U(1) port consists of $\sim 5400$ lines of generic block-sparse machinery
copied verbatim from \texttt{KagomeDSF/U1/}, plus $\sim 300$ lines of native
XXZ-specific code (\texttt{build\_xxz\_mpo\_charged} and helpers).  The MPO
charge structure is determined by the existing $\chi_{\rm MPO}=5$ FSM with
bond charges $(0,+2,-2,0,0)$, and the construction reduces mechanically to
``build dense $W$, project with \texttt{from\_dense}''.  No new algorithms
are needed; \texttt{dmrg\_sweep!}, \texttt{tdvp\_sweep!}, and
\texttt{tdvp\_sweep\_one\_site!} carry over with identical sweep mechanics.
Validation pairs every new path against both the existing dense baseline
and ED on small $N$.  Estimated effort: $\sim 10$ working days from
skeleton to a working end-to-end quench protocol script at $N=16$ outside
ED reach.

\end{document}
